\begin{nota}[Equivariant residue realization]\label{bipo:gdv}
Write $\ggk^{\widehat\mu,\mathrm{raw}}\var_{\C}$ for the ring with good finite
$\widehat\mu$-actions and the equivariant scissor relations.
Here \emph{good} means that the action factors through some $\mu_N$
and every orbit is contained in an affine open subset.
Write $\gmv$ for its quotient by the ideal generated by
\[
 [X\times(\A^r,\rho)]-[X\times\A^r],
\]
for every good equivariant $X$ and every finite linear action
$\rho$; the action on the second affine factor on the right is
trivial. Goodness and this additional relation are separate
conditions.

There is a canonical isomorphism
\[
 \Theta:\sggk\RES\xrightarrow{\ \sim\ }\gmv .
\]
We construct it at a finite Kummer level and prove both
directions of the comparison below.
\end{nota}

\begin{proof}
\emph{Finite-level coordinates and the raw comparison.}
For $q\in\Q$, let $L_q=q^\sharp\sqcup\{0_q\}$ be the
one-dimensional residue vector space whose nonzero part is
$q^\sharp=\vrv^{-1}(q)$. The notation $0_q$ records its zero
in that fixed vector space, encoded using the existing RV
zero/infinity element. It introduces no second RV zero.

For $U\in\RES$, split its finite value support and its zero
coordinate positions. Each piece can be regarded as a definable
subset of $\prod_{j=1}^r L_{a_j/N}$ for a common $N>0$ and
integers $a_j$. Choose one root $\tau_N^N=t$ and put
$r_N=\rv(\tau_N)$. The coordinate map
\[
 g_{r_N}(u_1,\ldots,u_r)
   =(u_1/r_N^{a_1},\ldots,u_r/r_N^{a_r})
\]
takes this product to $\C^r$, with $0_q$ mapped to $0$.
Residue-field purity makes $X_U=g_{r_N}(U)$ constructible.
The action of $\mu_N$ on its coordinates is
\[
 (x_j)_j\longmapsto(\zeta^{a_j}x_j)_j .
\]
The set $X_U$ is invariant because $U$ is defined over $F$.
Replacing $r_N$ by $\eta r_N$ multiplies the $j$th coordinate
by $\eta^{-a_j}$, an equivariant algebraic isomorphism.
Enlarging $N$ only inflates the action along
$\mu_M\to\mu_N$ and changes coordinates by such constants.
Thus the resulting equivariant constructible class is independent
of both choices. Products use a common Kummer level, so the
assignment respects products as well as scissors.

We now prove that every equivariant constructible set and
bijection is obtained in this way. A good finite action admits a
finite invariant affine stratification: intersect the translates
of an affine neighborhood of each orbit, take a finite such
cover, and successively remove the preceding opens. These
intersections are affine because the varieties are separated.
For an affine $\mu_N$-variety $X$, take algebra generators of
$\C[X]$, span their finite orbits, and diagonalize that finite
dimensional cyclic representation. Its eigenvectors still
generate the algebra. They give an equivariant closed embedding
\[
 X\hookrightarrow\A^r,\qquad
 \zeta\cdot(x_j)_j=(\zeta^{a_j}x_j)_j
\]
for suitable integers $a_j$.

Let
\[
 T_N=\{r\in\RV:r^N=\rv(t)\}.
\]
For any invariant constructible
$X\subseteq\A^r$ with these coordinate weights, define
\[
 U_X=\left\{u\in\prod_jL_{a_j/N}:
       \exists r\in T_N\quad g_r(u)\in X\right\}.
\]
This is an original-language $F$-definable RES set.
Every two elements of $T_N$ differ by an element of $\mu_N$.
Invariance of $X$ therefore implies
$g_r(U_X)=X$ for every $r\in T_N$.
This proves surjectivity without adding a root to the base.

The same construction applies to the graph of an equivariant
constructible bijection $f:X\to Y$, using one common level:
\[
 G_f=\{(u,v):\exists r\in T_N,\
              (g_r(u),g_r(v))\in\operatorname{Graph}(f)\}.
\]
Equivariance makes the condition independent of $r$, so $G_f$
is the graph of an $F$-definable bijection.
Conversely, framing the graph of an $F$-definable RES bijection
gives an equivariant constructible bijection.
In characteristic zero such a bijection is piecewise an
algebraic isomorphism: its generic graph over each irreducible
piece has degree one, hence a rational inverse on a dense open;
remove those opens and induct on dimension. Intersecting the
finitely many translates makes the partition invariant.
Consequently these graph constructions also respect the usual
equivariant scissor presentation. They show independence of the
chosen affine embeddings and stratifications and give inverse
maps on classes. We have proved an isomorphism
\[
 \Theta_{\mathrm{raw}}:\ggk\RES\xrightarrow{\ \sim\ }\ggk^{\widehat\mu,\mathrm{raw}}\var_{\C} .
\]

\emph{The quotient.}
For $q=a/N$, the space $L_q$ is sent to the affine line with
character $\rho_a(\zeta)x=\zeta^a x$. Its zero is a fixed point.
Hence
\[
 \Theta_{\mathrm{raw}}
 \bigl([q^\sharp]-[\G_m]\bigr)
 =
 [\A^1,\rho_a]-[\A^1,\mathrm{triv}].
\]
Every finite character of $\widehat\mu$ has this form.
Every finite linear cyclic action is a sum of characters, so its
class is a product of these line classes. By expanding the
difference of two such products, all the linear-action relations,
also after multiplication by an arbitrary $[X]$, are generated
by the displayed line differences. The two quotient ideals
therefore correspond exactly. This proves the asserted
isomorphism $\Theta$.

For later computations, let
$h_{M_i}(U_i)=V_i\times I_i^\sharp$ be a monomial decomposition
as in Lemma~5.12, with $\ell_i$ free RV coordinates.
The original-language construction of $\bb E_b$ in Section~4 gives
\begin{equation}\label{twist:C}
 (\Theta\circ\bb E_b)([\bm U])
   =\sum_i\chi_b(I_i)\,\Theta([V_i])[\G_m]^{\ell_i}.
\end{equation}
The finite-to-one label and its degree are retained as in the
product construction of Section~3; the displayed target forgets
that grading. The formula is an evaluation of the already defined
retraction, so it is independent of the decomposition.
\end{proof}

\begin{nota}\label{mov:vol}
The motivic volume operator is
\[
 \Vol=\Theta\circ\bb E_b\circ\int:
 \ggk\VF_*\longrightarrow\gmv .
\]
Here $\bb E_b$ is applied after the quotient by $(\bm P-1)$,
on which it is already defined by Section~4.
\end{nota}

Let $\mathcal A\subseteq\Gamma^n\times\VF^m\times\RV_\infty^s$
be an $F$-definable family whose fibers $A_\gamma$ are objects
of $\VF_*$. All families below are first interpreted in a
sufficiently saturated elementary extension $\mathbb U$.
They are restricted to $\widetilde{\C}$ only after uniform
representatives and their witnesses have been obtained.

For every $\gamma\in\Gamma(\mathbb U)^n$,
the base generated by $F$ and $\gamma$ is
$(\VF,\Gamma)$-generated. Apply Theorem~\ref{main:prop:HK}{} over that
base to choose RV representatives $\bm U_\gamma$ and actual
definable bijections
\[
 H_\gamma:A_\gamma\xrightarrow{\ \sim\ }\mathbb L\bm U_\gamma.
\]
These are equalities of positive classes, before groupification.
We retain the grading, the finite-to-one structure maps and
their finite-fiber bounds in the representative data.

Here is the precise compactness argument. Fix formulas for one
such representative and its bijection, with only the indicated
parameter $\gamma$ in addition to parameters from $F$.
For these formulas, being valid representative data is a
first-order condition on $\gamma$: it asserts the chosen sort
domains, totality of the structure maps, the specified finite-fiber
bounds, and that the graph of $H_\gamma$ defines a bijection with
the indicated domain and codomain.
The conditions from all choices cover the entire parameter
space in the saturated extension. Compactness gives a finite
subcover. Its finite Boolean refinement is an $F$-definable
partition on each piece of which the RV arities, label degrees,
representatives and bijection witnesses are given by fixed
formulas. Empty fibers are included with empty representatives.

Using tags for the finitely many pieces and degrees, write
$\mathcal U=\{(\gamma,u):u\in U_\gamma\}$ for the resulting
definable RV family; the structure maps remain part of its data.
In the inflated families, distinguish
\[
 A^{\mathrm{inf}}
   =\{(x,a):(\vv(x),a)\in\mathcal A\},
 \qquad
 U^{\mathrm{inf}}
   =\{(z,u):(\vrv(z),u)\in\mathcal U\},
\]
where $x\in(\VF^\times)^n$ and $z\in(\RV^\times)^n$.
Thus the first construction uses the VF value fibers and the
second uses the RV value fibers. The uniform bijections
$H_\gamma$ justify replacement of the inflated total VF space
by the corresponding family of lifts. Restricting all this
data to the Puiseux field yields the required representatives
at rational parameters.


% BEGIN AUTHOR RESOLUTION THETA
{\normalfont\small\color{ResolutionGreen}
\noindent\textbf{Author resolution (finite-level comparison).}
The author approved the original-language construction of $\Theta$ and its proof on November 13, 2026. The inverse uses finite Kummer root sets and equivariant graphs, and the proof identifies the exact quotient ideals. The comparison is proved here rather than assumed from the cited source.\par}
% END AUTHOR RESOLUTION THETA

% BEGIN REFEREE INSERTION R083
\ResolvedRefereeRemark{083}{Clarification; revised}{Uniform representatives and compactness}{%
\textit{Original source line 1825.}\quad Make this compactness step precise by first interpreting the family in a saturated elementary extension. Apply Theorem~\ref{main:prop:HK} over $\mathbb S\cup\{\gamma\}$ there, uniformize the representatives together with their finite bijection witnesses, and then restrict to the Puiseux model. Compactness cannot be applied just to an arbitrary list of representatives chosen at rational points. Also distinguish the valued-field fibers $\vv^{-1}(\gamma)$ used in $A$ from the leading-term fibers $\vrv^{-1}(\gamma)$ used in $U$.
}{The author approved the replacement on November 13, 2026. Representatives, grading and structure maps, finite-fiber bounds, and actual lift bijections are uniformized together in a saturated elementary extension over the generated $(\VF,\Gamma)$ base. Only after the finite definable partition is obtained is the data restricted to Puiseux rational parameters. VF and RV value fibers are distinguished explicitly.}
% END REFEREE INSERTION R083

